\textbf{Phần 1. Trắc nghiệm lựa chọn}
\setcounterpageref{ex}{0}

\begin{ex} %[2D13V66]
Một hộp chỉ đựng toàn bóng màu đỏ. Nếu nhắm mắt lấy ra một quả bóng thì biến cố nào sau đây là 'Chắc chắn' xảy ra?
\choice
{Lấy được quả bóng màu xanh}
{\True Lấy được quả bóng màu đỏ}
{Lấy được quả bóng màu vàng}
{Không lấy được quả nào}
\loigiai{
Hộp toàn bóng đỏ nên chắc chắn lấy ra quả bóng đỏ.
\textbf{Đáp án của câu hỏi là B}
}
\end{ex}

\begin{ex} %[2D13V66]
Trong túi có 3 viên bi xanh và 2 viên bi đỏ. Lấy ngẫu nhiên một viên bi. Khẳng định nào sau đây là ĐÚNG?
\choice
{Chắc chắn lấy được viên bi màu vàng}
{Không thể lấy được viên bi màu xanh}
{Chắc chắn lấy được bi đỏ}
{\True Có thể lấy được viên bi màu đỏ}
\loigiai{
Vì trong túi có cả bi xanh và đỏ nên có thể lấy được bi đỏ (hoặc xanh).
\textbf{Đáp án của câu hỏi là D}
}
\end{ex}

\begin{ex} %[2D13V66]
Gieo một con xúc xắc có các mặt từ 1 đến 6 chấm. Biến cố 'Xuất hiện mặt 8 chấm' là sự kiện
\choice
{\True Không thể xảy ra}
{Chắc chắn xảy ra}
{Có thể xảy ra}
{Luôn luôn xảy ra}
\loigiai{
Xúc xắc chỉ có từ 1 đến 6 chấm nên không thể xuất hiện mặt 8 chấm.
\textbf{Đáp án của câu hỏi là A}
}
\end{ex}

\begin{ex} %[2D13V66]
Hôm nay là ngày thứ Hai. Biến cố 'Ngày mai là ngày thứ Ba' là sự kiện
\choice
{Có thể xảy ra}
{Không thể xảy ra}
{\True Chắc chắn xảy ra}
{Không xác định}
\loigiai{
Thứ Hai luôn luôn liền kề trước thứ Ba nên chắc chắn xảy ra.
\textbf{Đáp án của câu hỏi là C}
}
\end{ex}

\begin{ex} %[2D13V66]
Một hộp có 5 chiếc bút chì màu đen. Bạn An thò tay vào hộp lấy ra 1 chiếc bút. Bạn An
\choice
{Có thể lấy được chiếc bút chì màu đỏ}
{\True Không thể lấy được chiếc bút chì màu đỏ}
{Chắc chắn lấy được bút chì màu đỏ}
{Chắc chắn lấy được tẩy}
\loigiai{
Trong hộp không có bút chì đỏ nên không thể lấy được.
\textbf{Đáp án của câu hỏi là B}
}
\end{ex}

\begin{ex} %[2D13V66]
Gieo một đồng xu có hai mặt Sấp và Ngửa. Kết quả xuất hiện mặt Ngửa là sự kiện
\choice
{Chắc chắn xảy ra}
{Không thể xảy ra}
{Luôn luôn sai}
{\True Có thể xảy ra}
\loigiai{
Đồng xu có thể rơi vào mặt sấp hoặc mặt ngửa.
\textbf{Đáp án của câu hỏi là D}
}
\end{ex}

\begin{ex} %[2D13V66]
Mặt trời mọc ở hướng Đông vào mỗi buổi sáng là sự kiện
\choice
{\True Chắc chắn xảy ra}
{Có thể xảy ra}
{Không thể xảy ra}
{Tùy từng ngày}
\loigiai{
Quy luật tự nhiên: Mặt trời luôn mọc ở hướng Đông nên chắc chắn xảy ra.
\textbf{Đáp án của câu hỏi là A}
}
\end{ex}

\begin{ex} %[2D13V66]
Một tuần có 8 ngày là sự kiện
\choice
{Có thể xảy ra}
{Chắc chắn xảy ra}
{\True Không thể xảy ra}
{Xảy ra vào năm nhuận}
\loigiai{
Một tuần luôn chỉ có đúng 7 ngày, do đó một tuần có 8 ngày là không thể.
\textbf{Đáp án của câu hỏi là C}
}
\end{ex}

\begin{ex} %[2D13V66]
Trong lọ có 4 viên kẹo cam và 4 viên kẹo dâu. Mai lấy ra 5 viên kẹo. Khẳng định nào sau đây là 'Chắc chắn'?
\choice
{Cả 5 viên kẹo đều là kẹo cam}
{\True Trong 5 viên kẹo lấy ra có ít nhất 1 viên kẹo cam}
{Cả 5 viên kẹo đều là kẹo dâu}
{Không có viên kẹo dâu nào}
\loigiai{
Vì chỉ có 4 viên dâu nên lấy 5 viên chắc chắn phải có ít nhất 1 viên kẹo cam.
\textbf{Đáp án của câu hỏi là B}
}
\end{ex}

\begin{ex} %[2D13V66]
Một em bé mới sinh biết đi xe đạp ngay lập tức là sự kiện
\choice
{Có thể xảy ra}
{Chắc chắn xảy ra}
{Thỉnh thoảng xảy ra}
{\True Không thể xảy ra}
\loigiai{
Em bé mới sinh không thể biết đi xe đạp.
\textbf{Đáp án của câu hỏi là D}
}
\end{ex}

\begin{ex} %[2D13V66]
Khi quay một chiếc kim trên vòng tròn chia làm hai nửa: nửa đỏ và nửa xanh. Kim dừng lại ở phần màu đỏ là sự kiện
\choice
{\True Có thể xảy ra}
{Chắc chắn xảy ra}
{Không thể xảy ra}
{Không bao giờ xảy ra}
\loigiai{
Kim có thể dừng ở màu đỏ hoặc màu xanh.
\textbf{Đáp án của câu hỏi là A}
}
\end{ex}

\begin{ex} %[2D13V66]
Tổng số chấm khi gieo hai con xúc xắc lớn hơn 1 là sự kiện
\choice
{Có thể xảy ra}
{Không thể xảy ra}
{\True Chắc chắn xảy ra}
{Tùy lúc}
\loigiai{
Mỗi con xúc xắc ít nhất 1 chấm, hai con ít nhất $1 + 1 = 2 > 1$, nên chắc chắn lớn hơn 1.
\textbf{Đáp án của câu hỏi là C}
}
\end{ex}

\textbf{Phần 2. Trắc nghiệm đúng sai}
\setcounterpageref{ex}{0}

\begin{ex} %[2D13V66]
Xét tính đúng sai của các nhận định về khả năng xảy ra của sự kiện:
\choiceTF
{\True Một sự kiện chắc chắn xảy ra nghĩa là nó luôn luôn diễn ra trong mọi trường hợp.}
{Một sự kiện có thể xảy ra nghĩa là nó không bao giờ xảy ra được.}
{\True Sự kiện con gà đẻ ra quả trứng vịt là sự kiện không thể xảy ra.}
{\True Lấy 1 viên kẹo từ túi có 10 viên kẹo sô-cô-la chắc chắn được viên kẹo sô-cô-la.}
\loigiai{
\begin{itemize}
    \item Ý a) Đúng: Đúng định nghĩa.
    \item Ý b) Sai: Có thể nghĩa là có khả năng xảy ra. Sai.
    \item Ý c) Đúng: Đúng thực tế sinh học.
    \item Ý d) Đúng: Đúng vì túi chỉ có kẹo sô-cô-la.
\end{itemize}
\begin{itemize}
    \item a) \textbf{Đúng}.
    \item b) \textbf{Sai}.
    \item c) \textbf{Đúng}.
    \item d) \textbf{Đúng}.
\end{itemize}
}
\end{ex}

\begin{ex} %[2D13V66]
Trong một hộp có 3 quả bóng xanh và 3 quả bóng đỏ:
\choiceTF
{\True Lấy ra 1 quả bóng thì có thể được bóng xanh.}
{\True Lấy ra 1 quả bóng thì không thể được bóng vàng.}
{Lấy ra 2 quả bóng thì chắc chắn cả 2 quả đều là bóng xanh.}
{\True Lấy ra 4 quả bóng thì chắc chắn có ít nhất 1 quả bóng màu xanh.}
\loigiai{
\begin{itemize}
    \item Ý a) Đúng: Đúng.
    \item Ý b) Đúng: Trong hộp không có bóng vàng. Đúng.
    \item Ý c) Sai: Có thể lấy được 1 xanh 1 đỏ hoặc 2 đỏ. Sai.
    \item Ý d) Đúng: Vì chỉ có 3 bóng đỏ nên 4 quả chắc chắn có ít nhất 1 xanh. Đúng.
\end{itemize}
\begin{itemize}
    \item a) \textbf{Đúng}.
    \item b) \textbf{Đúng}.
    \item c) \textbf{Sai}.
    \item d) \textbf{Đúng}.
\end{itemize}
}
\end{ex}

\textbf{Phần 3. Trả lời ngắn}
\setcounterpageref{ex}{0}

\begin{ex} %[2D13V66]
Trong hộp có 5 viên bi đỏ và 1 viên bi xanh. Hỏi cần lấy ra ít nhất bao nhiêu viên bi để chắc chắn có 1 viên bi màu đỏ?
\shortans{2}
\loigiai{
\textbf{Đáp số:} $2$\\
Trường hợp xấu nhất lấy phải 1 viên bi xanh, lấy thêm 1 viên nữa chắc chắn là bi đỏ. Cần lấy $1 + 1 = 2$ viên.
}
\end{ex}

\begin{ex} %[2D13V66]
Một con xúc xắc có 6 mặt từ 1 đến 6 chấm. Có bao nhiêu mặt có số chấm là số chẵn?
\shortans{3}
\loigiai{
\textbf{Đáp số:} $3$\\
Các mặt chẵn là 2, 4, 6. Có 3 mặt.
}
\end{ex}

\begin{ex} %[2D13V66]
Trong túi có 10 chiếc kẹo gồm hai loại: kẹo cam và kẹo dâu. Biết chắc chắn lấy được ít nhất 1 chiếc kẹo dâu khi lấy 6 chiếc kẹo bất kì. Có bao nhiêu chiếc kẹo cam trong túi?
\shortans{5}
\loigiai{
\textbf{Đáp số:} $5$\\
Để chắc chắn có kẹo dâu khi lấy 6 chiếc thì số kẹo cam nhiều nhất là 5 chiếc.
}
\end{ex}

\begin{ex} %[2D13V66]
Gieo 1 con xúc xắc, có mấy khả năng xuất hiện mặt có số chấm lớn hơn 4?
\shortans{2}
\loigiai{
\textbf{Đáp số:} $2$\\
Mặt 5 và mặt 6, có 2 khả năng.
}
\end{ex}

\begin{ex} %[2D13V66]
Trong hộp có 4 chiếc bút chì vàng. An thò tay lấy 2 chiếc. Có bao nhiêu chiếc bút chì màu khác vàng được lấy ra?
\shortans{0}
\loigiai{
\textbf{Đáp số:} $0$\\
Trong hộp chỉ có bút vàng nên lấy được 0 chiếc màu khác.
}
\end{ex}

\begin{ex} %[2D13V66]
Một tuần lễ chắc chắn có mấy ngày?
\shortans{7}
\loigiai{
\textbf{Đáp số:} $7$\\
Một tuần lễ chắc chắn có đúng 7 ngày.
}
\end{ex}
